% Options for packages loaded elsewhere
\PassOptionsToPackage{unicode}{hyperref}
\PassOptionsToPackage{hyphens}{url}
\documentclass[
]{article}
\usepackage{xcolor}
\usepackage{amsmath,amssymb}
\usepackage{graphicx}
\providecommand{\llbracket}{\mathopen{[\![}}
\providecommand{\rrbracket}{\mathclose{]\!]}}
\providecommand{\coloneqq}{:=}
\setcounter{secnumdepth}{-\maxdimen} % remove section numbering
\usepackage{iftex}
\ifPDFTeX
  \usepackage[T1]{fontenc}
  \usepackage[utf8]{inputenc}
  \usepackage{textcomp} % provide euro and other symbols
\else % if luatex or xetex
  \usepackage{unicode-math} % this also loads fontspec
  \defaultfontfeatures{Scale=MatchLowercase}
  \defaultfontfeatures[\rmfamily]{Ligatures=TeX,Scale=1}
\fi
\usepackage{lmodern}
\ifPDFTeX\else
  % xetex/luatex font selection
\fi
% Use upquote if available, for straight quotes in verbatim environments
\IfFileExists{upquote.sty}{\usepackage{upquote}}{}
\IfFileExists{microtype.sty}{% use microtype if available
  \usepackage[]{microtype}
  \UseMicrotypeSet[protrusion]{basicmath} % disable protrusion for tt fonts
}{}
\makeatletter
\@ifundefined{KOMAClassName}{% if non-KOMA class
  \IfFileExists{parskip.sty}{%
    \usepackage{parskip}
  }{% else
    \setlength{\parindent}{0pt}
    \setlength{\parskip}{6pt plus 2pt minus 1pt}}
}{% if KOMA class
  \KOMAoptions{parskip=half}}
\makeatother
\usepackage{color}
\usepackage{fancyvrb}
\newcommand{\VerbBar}{|}
\newcommand{\VERB}{\Verb[commandchars=\\\{\}]}
\DefineVerbatimEnvironment{Highlighting}{Verbatim}{commandchars=\\\{\}}
% Add ',fontsize=\small' for more characters per line
\newenvironment{Shaded}{}{}
\newcommand{\AlertTok}[1]{\textcolor[rgb]{1.00,0.00,0.00}{\textbf{#1}}}
\newcommand{\AnnotationTok}[1]{\textcolor[rgb]{0.38,0.63,0.69}{\textbf{\textit{#1}}}}
\newcommand{\AttributeTok}[1]{\textcolor[rgb]{0.49,0.56,0.16}{#1}}
\newcommand{\BaseNTok}[1]{\textcolor[rgb]{0.25,0.63,0.44}{#1}}
\newcommand{\BuiltInTok}[1]{\textcolor[rgb]{0.00,0.50,0.00}{#1}}
\newcommand{\CharTok}[1]{\textcolor[rgb]{0.25,0.44,0.63}{#1}}
\newcommand{\CommentTok}[1]{\textcolor[rgb]{0.38,0.63,0.69}{\textit{#1}}}
\newcommand{\CommentVarTok}[1]{\textcolor[rgb]{0.38,0.63,0.69}{\textbf{\textit{#1}}}}
\newcommand{\ConstantTok}[1]{\textcolor[rgb]{0.53,0.00,0.00}{#1}}
\newcommand{\ControlFlowTok}[1]{\textcolor[rgb]{0.00,0.44,0.13}{\textbf{#1}}}
\newcommand{\DataTypeTok}[1]{\textcolor[rgb]{0.56,0.13,0.00}{#1}}
\newcommand{\DecValTok}[1]{\textcolor[rgb]{0.25,0.63,0.44}{#1}}
\newcommand{\DocumentationTok}[1]{\textcolor[rgb]{0.73,0.13,0.13}{\textit{#1}}}
\newcommand{\ErrorTok}[1]{\textcolor[rgb]{1.00,0.00,0.00}{\textbf{#1}}}
\newcommand{\ExtensionTok}[1]{#1}
\newcommand{\FloatTok}[1]{\textcolor[rgb]{0.25,0.63,0.44}{#1}}
\newcommand{\FunctionTok}[1]{\textcolor[rgb]{0.02,0.16,0.49}{#1}}
\newcommand{\ImportTok}[1]{\textcolor[rgb]{0.00,0.50,0.00}{\textbf{#1}}}
\newcommand{\InformationTok}[1]{\textcolor[rgb]{0.38,0.63,0.69}{\textbf{\textit{#1}}}}
\newcommand{\KeywordTok}[1]{\textcolor[rgb]{0.00,0.44,0.13}{\textbf{#1}}}
\newcommand{\NormalTok}[1]{#1}
\newcommand{\OperatorTok}[1]{\textcolor[rgb]{0.40,0.40,0.40}{#1}}
\newcommand{\OtherTok}[1]{\textcolor[rgb]{0.00,0.44,0.13}{#1}}
\newcommand{\PreprocessorTok}[1]{\textcolor[rgb]{0.74,0.48,0.00}{#1}}
\newcommand{\RegionMarkerTok}[1]{#1}
\newcommand{\SpecialCharTok}[1]{\textcolor[rgb]{0.25,0.44,0.63}{#1}}
\newcommand{\SpecialStringTok}[1]{\textcolor[rgb]{0.73,0.40,0.53}{#1}}
\newcommand{\StringTok}[1]{\textcolor[rgb]{0.25,0.44,0.63}{#1}}
\newcommand{\VariableTok}[1]{\textcolor[rgb]{0.10,0.09,0.49}{#1}}
\newcommand{\VerbatimStringTok}[1]{\textcolor[rgb]{0.25,0.44,0.63}{#1}}
\newcommand{\WarningTok}[1]{\textcolor[rgb]{0.38,0.63,0.69}{\textbf{\textit{#1}}}}
\setlength{\emergencystretch}{3em} % prevent overfull lines
\providecommand{\tightlist}{%
  \setlength{\itemsep}{0pt}\setlength{\parskip}{0pt}}
\usepackage{bookmark}
\IfFileExists{xurl.sty}{\usepackage{xurl}}{} % add URL line breaks if available
\urlstyle{same}
\hypersetup{
  hidelinks,
  pdfcreator={LaTeX via pandoc}}

\author{}
\date{}

\begin{document}

\section{Formalization (Pure Automata
Version)}\label{formalization-pure-automata-version}

\subsection{1. Scope and Goal}\label{scope-and-goal}

Scope:

\begin{itemize}
\tightlist
\item
  synchronous Kairos programs modeled in abstract relational form,
\item
  temporal contracts in the safety LTL fragment,
\item
  translation of LTL contracts into local first-order obligations over
  finite histories, proved with standard deductive verification
  techniques.
\end{itemize}

Goal:

Given a node \(N\) (program semantics) and a specification
\(\Phi=(A,G)\), generate and prove local first-order obligations over
finite histories such that:

\[
\forall u,\;
u \models A
\implies
run(N,u) \models G.
\]

where \(u\) is an input stream (values read at each logical tick), and
\(run(N,u)\) is the observable trace produced by \(N\) on \(u\). In
short, every admissible input stream (satisfying \texttt{A}) induces
only observable behaviors satisfying \texttt{G}.

\subsection{2. Logic and Contracts}\label{logic-and-contracts}

\subsubsection{2.1 Streams and Interpretation
Domains}\label{streams-and-interpretation-domains}

Let \(I\) be the set of input variables and \(O\) the set of output
variables.

For any variable set \(V\), let \(\mathrm{Val}(V)\) be the set of
valuations of \(V\), and \(\mathrm{Str}(V)=\mathrm{Val}(V)^\omega\) the
set of infinite streams of valuations.

No finiteness assumption is made on value domains: \(\mathrm{Val}(V)\)
may be infinite (for example with integer-valued variables).

We use input streams \(u \in \mathrm{Str}(I)\) and observable streams
\(w \in \mathrm{Str}(I \cup O)\) as semantic domains.

Assumptions constrain input behavior, so they are interpreted on
\(\mathrm{Str}(I)\); guarantees constrain observable behavior, so they
are interpreted on \(\mathrm{Str}(I \cup O)\). For a node \(N\), Section
3 will define how each input stream \(u\) induces one observable stream
\(run(N,u)\in\mathrm{Str}(I\cup O)\).

We write \(\mathrm{Hist}(V)\) for the set of finite backward histories
over \(V\), and \(\mathrm{HFO}(V)\) for first-order predicates over
\(\mathrm{Hist}(V)\). At time \(i\), \(h_i^V\in\mathrm{Hist}(V)\)
denotes the current finite history prefix used to interpret state
predicates. Hence \(A\) is an LTL formula whose state propositions are
in \(\mathrm{HFO}(I)\), and \(G\) is an LTL formula whose state
propositions are in \(\mathrm{HFO}(I\cup O)\).

\subsubsection{2.2 Formula Fragment and
Satisfaction}\label{formula-fragment-and-satisfaction}

We use formulas built from temporal operators \texttt{X} (next) and
\texttt{W} (weak until), with standard Boolean connectives
(\(\land,\lor,\neg\)); \(\Rightarrow\) and \(\Leftrightarrow\) are
abbreviations.

In this document, contract formulas \(A\) and \(G\) are assumed to
satisfy the safety-oriented positivity condition: \texttt{W} occurs only
positively (never under negation). This condition is checked
syntactically and is used later to build deterministic safety automata.

Grammar: \[
\varphi ::= p \mid \neg\varphi \mid \varphi \land \varphi \mid \varphi \lor \varphi \mid X\varphi \mid \varphi \,W\, \varphi,
\qquad p\in\mathrm{HFO}(V).
\]

We also use the standard global operator \(\mathbf{G}\) as an
abbreviation: \[
\mathbf{G}\varphi \;\equiv\; \varphi\,W\,false.
\]

Satisfaction for \(\pi\in\mathrm{Str}(V)\), \(i\in\mathbb N\): \[
\begin{array}{l}
\pi,i \models p \iff p(h_i^V) \text{ where } h_i^V \in \mathrm{Hist}(V) \text{ is the finite history at time } i\\
\pi,i \models \neg\varphi \iff \neg(\pi,i\models\varphi)\\
\pi,i \models \varphi\land\psi \iff (\pi,i\models\varphi)\land(\pi,i\models\psi)\\
\pi,i \models \varphi\lor\psi \iff (\pi,i\models\varphi)\lor(\pi,i\models\psi)\\
\pi,i \models X\varphi \iff \pi,i+1\models\varphi\\
\pi,i \models \varphi W \psi \iff
\Bigl(\forall j\ge i,\ \pi,j\models\varphi\Bigr)\ \lor
\Bigl(\exists k\ge i,\ \pi,k\models\psi\ \land\ \forall j\in[i,k),\ \pi,j\models\varphi\Bigr)\\
\pi,i \models \mathbf{G}\varphi \iff \forall j\ge i,\ \pi,j\models\varphi\\
\quad\text{and}\quad
\pi\models\varphi \iff \pi,0\models\varphi
\end{array}
\]

\subsubsection{2.3 Safety Property
Assumption}\label{safety-property-assumption}

From the positivity condition above (with Boolean abbreviations
expanded), contracts are taken in the safety fragment: \[
A \in \mathrm{LTL}_{safe}(\mathrm{HFO}(I)),
\qquad
G \in \mathrm{LTL}_{safe}(\mathrm{HFO}(I\cup O)).
\]

Formally, \(\varphi\) is safety iff: \[
\forall \pi \in \mathrm{Str}(V),\;
\pi \not\models \varphi
\Rightarrow
\exists n\in\mathbb N,\;
\forall \pi' \in \mathrm{Str}(V),\;
\pi[0..n]\cdot \pi' \not\models \varphi.
\]

Intuition: every violation has a finite bad prefix.

The automata-level view used by the method is introduced in Section 5.

\subsection{3. Abstract Kairos Model}\label{abstract-kairos-model}

This abstraction follows the synchronous reactive tradition and is close
to a symbolic Mealy/Lustre-style model.

Kairos is viewed as a synchronous transition language: at each logical
tick, a node reads inputs, evaluates control conditions (\texttt{when}),
and produces outputs and next-state values through guarded updates.

The objective is a relational, proof-oriented semantics of Kairos
programs so temporal contracts can be reduced to local first-order proof
obligations over transitions.

We separate the program from its specification. This abstract model is
used to simplify the presentation; the tool does not explicitly
translate programs into this model as a standalone intermediate
artifact.

\subsubsection{3.1 Syntax}\label{syntax}

\paragraph{3.1.1 Node and Specification}\label{node-and-specification}

A node is represented as \(N=(\mathcal S,s_0,\mathcal T,X,I,O)\), where
\(\mathcal S\) is the finite control-state set, \(s_0\in\mathcal S\) is
the initial control state, \(X\) is the variable set, and
\(I,O\subseteq X\) are the input and output variables with
\(I\cap O=\varnothing\). Internal variables are
\(L=X\setminus(I\cup O)\), so \(X=I\uplus O\uplus L\).

We write transitions with the arrow notation
\(s \xrightarrow{\gamma/\rho} s'\), equivalently
\((s,\gamma,\rho,s')\in\mathcal T\).

Let \(\Sigma=\mathrm{Val}(X)\) be the valuation domain over program
variables. We use \(\Gamma=(\Sigma\to\mathbb B)\) for guards and
\(R=(\Sigma\times\Sigma\to\mathbb B)\) for update relations, so
\(\mathcal T\subseteq\mathcal S\times\Gamma\times R\times\mathcal S\).
Intuitively, \(\gamma\) constrains the current step context and \(\rho\)
constrains the pair (current valuation, next valuation). Guards are
interpreted on the current finite history \(h_i^X\).

Update relations are intentionally open-world in this abstraction: if a
next-state variable is not constrained by \(\rho\), it remains
unconstrained.

Initialization convention: before the first logical step, internal and
output variables may be undefined. Here, ``undefined'' means ``has no
value yet'', not ``can take an arbitrary value''. The \texttt{Init}
control phase is used to establish defined values before regular-state
reads.

The specification associated with a node is \(\Phi=(A,G,Inv)\), where
\texttt{A} is a temporal assumption on input traces, \texttt{G} is a
temporal guarantee on observable traces, and
\(Inv:\mathcal S\to\mathrm{HFO}(X)\) maps each control state to a
first-order state invariant over finite histories.

Types of \texttt{A} and \texttt{G} are those introduced in Section 2.
Correctness is established for the pair \((N,\Phi)\).

\paragraph{3.1.2 Well-Formedness (Deterministic
Profile)}\label{well-formedness-deterministic-profile}

We target deterministic, non-blocking programs. Since the transition
model is relational, this is not automatic: from one configuration
\((s,\sigma)\), several control branches may be enabled, or one branch
may admit several next valuations. We therefore impose two high-level
constraints.

In this document, these constraints are assumptions of the model and are
considered true by construction of well-formed Kairos programs.

\texttt{Det-Step} (total functional step): \[
\forall s\in\mathcal S,\forall \sigma\in\Sigma,\exists!\,s'\in\mathcal S,\exists!\,\sigma'\in\Sigma,\;
(s,\sigma)\to(s',\sigma').
\]

\texttt{Input-Frame} (inputs are read-only across one step): \[
\forall (s,\sigma)\to(s',\sigma'),\;
\sigma\!\upharpoonright_I=\sigma'\!\upharpoonright_I.
\]

\paragraph{3.1.3 Abstract Examples}\label{abstract-examples}

For readability, we present only the node part in these examples and
omit the associated specification \(\Phi\) at this stage.

\texttt{delay\_int} (one-tick integer delay): \[
\begin{array}{l}
N_{delay}=(\{Init,Run\},Init,\mathcal T_{delay},X,\{x\},\{y\})\quad\text{with}\quad X=\{x,y,z\}\\
\mathcal T_{delay}=\{Init\xrightarrow{true\,/\,z'=x}Run,\;Run\xrightarrow{true\,/\,y'=z\land z'=x}Run\}
\end{array}
\]

\texttt{toggle} (alternating output): \[
\begin{array}{l}
N_{toggle}=(\{Init,Run\},Init,\mathcal T_{toggle},\{y\},\varnothing,\{y\})\\
\mathcal T_{toggle}=\{Init\xrightarrow{true\,/\,y'=0}Run,\;Run\xrightarrow{true\,/\,y'=1-y}Run\}
\end{array}
\]

\subsection{4. Node Semantics}\label{node-semantics}

\subsubsection{4.1 Step Semantics}\label{step-semantics}

The operational semantics of \(N\) is a transition relation on
configurations \((s,\sigma)\), where \(s \in \mathcal S\) and
\(\sigma \in \Sigma\). Intuitively, \((s,\sigma)\) moves to
\((s',\sigma')\) when there is a control transition
\((s,\gamma,\rho,s') \in \mathcal T\), \(\sigma\) satisfies \(\gamma\),
and \((\sigma,\sigma')\) satisfies \(\rho\).

Formally:

\[
(s,\sigma) \to (s',\sigma')
\iff
\exists \gamma,\rho.\;
(s,\gamma,\rho,s')\in\mathcal T
\;\land\;
\sigma\models\gamma
\;\land\;
(\sigma,\sigma')\models\rho
\]

This step relation is intra-instant: it computes the end-of-tick
valuation from the begin-of-tick valuation under the current input.

Under the \texttt{Det-Step} assumption from Section 3.1.2, this relation
is total and functional. We therefore write: \[
step_N:\mathcal S\times\Sigma\to\mathcal S\times\Sigma,\qquad
step_N(s,\sigma)=(s',\sigma')\iff (s,\sigma)\to(s',\sigma').
\]

\subsubsection{4.2 Stream Semantics}\label{stream-semantics}

An input stream is a sequence:

\[
u = u_0,u_1,u_2,\dots
\quad\text{with}\quad
u_i \in \{\sigma\!\upharpoonright_I \mid \sigma \in \Sigma\}.
\]

For a given input stream \(u\), an execution is a sequence of
begin-of-tick configurations \((s_i,\sigma_i)_{i\ge 0}\) such that, for
every \(i\), there exists an end-of-tick valuation
\(\sigma_i^+ \in \Sigma\) satisfying:

\[
\begin{array}{l}
\text{Step}(i):\ (s_i,\sigma_i)\to(s_{i+1},\sigma_i^+)\\
\text{SampleIn}(i):\ \sigma_i\!\upharpoonright_I = u_i\\
\text{Carry}(i):\ \sigma_{i+1}\!\upharpoonright_{X\setminus I} = \sigma_i^+\!\upharpoonright_{X\setminus I}\\
\text{SampleIn}(i+1):\ \sigma_{i+1}\!\upharpoonright_I = u_{i+1}
\end{array}
\]

The two \texttt{SampleIn} clauses encode input sampling at each tick.
\texttt{Carry} encodes state propagation: non-input variables at the
next tick start come from the previous tick end.

The observable trace induced by \(u\) is
\(run(N,u)=\tau_0,\tau_1,\tau_2,\dots\) over \(I\cup O\), where for
every \(i\): \[
\tau_i\!\upharpoonright_I=u_i
\qquad\text{and}\qquad
\tau_i\!\upharpoonright_O=\sigma_i^+\!\upharpoonright_O.
\]

Consistently with the initialization convention of Section 3, variables
may be undefined before the first step. In examples, when an initial
output is intentionally not constrained, we only comment the guaranteed
part of the output stream.

Example (\texttt{delay\_int}):

Using \(N_{delay}\) from Section 3.1.3 and input stream
\(x_0=3,\ x_1=5,\ x_2=2,\dots\), one execution is:

\[
\begin{array}{l}
(Init,\sigma_0)\to(Run,\sigma_0^+),\quad \sigma_0(x)=3,\ \sigma_0^+(z)=3\\
(Run,\sigma_1)\to(Run,\sigma_1^+),\quad \sigma_1(x)=5,\ \sigma_1(z)=3,\ \sigma_1^+(y)=3,\ \sigma_1^+(z)=5\\
(Run,\sigma_2)\to(Run,\sigma_2^+),\quad \sigma_2(x)=2,\ \sigma_2(z)=5,\ \sigma_2^+(y)=5,\ \sigma_2^+(z)=2
\end{array}
\]

Hence: \[
y_1 = x_0,\quad y_2 = x_1,\ \dots
\] so \(y\) is one tick behind \(x\) (after initialization). The value
\(y_0\) is not constrained by the \texttt{Init} transition in this
abstract example.

Example (\texttt{toggle}):

Using \(N_{toggle}\) from Section 3.1.3, one execution (begin-of-tick
valuations \(\sigma_i\), end-of-tick valuations \(\sigma_i^+\)) is:

\[
\begin{array}{l}
(Init,\sigma_0)\to(Run,\sigma_0^+),\quad \sigma_0^+(y)=0\\
(Run,\sigma_1)\to(Run,\sigma_1^+),\quad \sigma_1(y)=0,\ \sigma_1^+(y)=1\\
(Run,\sigma_2)\to(Run,\sigma_2^+),\quad \sigma_2(y)=1,\ \sigma_2^+(y)=0\\
(Run,\sigma_3)\to(Run,\sigma_3^+),\quad \sigma_3(y)=0,\ \sigma_3^+(y)=1
\end{array}
\]

so the observable output stream is \(0,1,0,1,\dots\).

Implementation correspondence: one logical tick in the implementation
matches one application of \texttt{step\_N}, plus input resampling and
state carry as above. The resulting stream \(run(N,u)\) is exactly the
semantic object consumed later by contract automata (Sections 5 to 7).

\subsubsection{4.3 Correctness w.r.t. the Temporal
Specification}\label{correctness-w.r.t.-the-temporal-specification}

Contract semantics is evaluated as follows: \texttt{A} is interpreted on
the input stream \(u\), and \texttt{G} is interpreted on the observable
stream \(run(N,u)\).

With \(\Phi=(A,G,Inv)\), we say that \(N\) is temporally correct w.r.t.
\(\Phi\) iff:

\[
\forall u,\;
u \models A
\implies
run(N,u) \models G.
\]

Intuitively: every admissible input stream (satisfying \texttt{A})
induces only observable behaviors that satisfy \texttt{G}.

\subsection{5. From Temporal Contracts to Safety
Automata}\label{from-temporal-contracts-to-safety-automata}

\subsubsection{5.1 Symbolic Safety
Automata}\label{symbolic-safety-automata}

Each safety formula \(\varphi\) over \(\mathrm{HFO}(V)\) can be compiled
into a deterministic symbolic safety automaton
\(\mathcal S_\varphi=(Q,q_0,\delta,bad)\), where: - \(Q\) is finite, -
\(q_0\in Q\) is initial, - \(\delta: Q \times \mathrm{Hist}(V) \to Q\)
is a total deterministic transition function, - \(bad\in Q\) is an
absorbing bad state.

Equivalently, for all \(q\in Q\) and \(h\in\mathrm{Hist}(V)\), there
exists exactly one \(q'\in Q\) such that \(\delta(q,h)=q'\).

For readability, we also use symbolic edge notation
\(q \xrightarrow{\alpha} q'\), where \(\alpha\in\mathrm{HFO}(V)\)
characterizes the histories \(h\) such that \(\delta(q,h)=q'\).

Given \(\pi\in\mathrm{Str}(V)\), we define its (unique) state sequence
\(r^\pi:\mathbb N\to Q\) by: \[
r^\pi(0)=q_0,\qquad r^\pi(i+1)=\delta(r^\pi(i),h_i^V).
\] The trace \(\pi\) is accepted iff
\(\forall i\in\mathbb N,\ r^\pi(i)\neq bad\).

Correctness of this automaton construction is the language equivalence:
\[
L(\varphi)=L(\mathcal S_\varphi)
=\{\pi\in\mathrm{Str}(V) \mid \pi \models \varphi\}.
\]

\subsubsection{\texorpdfstring{5.2 Contract Automata for \(A\) and
\(G\)}{5.2 Contract Automata for A and G}}\label{contract-automata-for-a-and-g}

Given \(\Phi=(A,G,Inv)\), we assume:

\begin{itemize}
\tightlist
\item
  \texttt{A} is compiled into a safety automaton
  \(\mathcal A=(Q_A,q^A_0,\Delta_A,bad_A)\),
\item
  \texttt{G} is compiled into a safety automaton
  \(\mathcal G=(Q_G,q^G_0,\Delta_G,bad_G)\).
\end{itemize}

The mapping \(Inv\) is used later to help the deductive proof method by
strengthening local proof contexts.

Assumptions on these automata:

\begin{itemize}
\tightlist
\item
  each symbolic edge guard is interpreted as a first-order predicate
  over finite histories;
\item
  automata are deterministic and complete (equivalently, they denote
  total transition functions on finite histories);
\item
  \(bad_A\) and \(bad_G\) are absorbing.
\end{itemize}

Standard correctness assumptions:

\begin{itemize}
\tightlist
\item
  automaton executions exactly follow the compiled LTL semantics,
\item
  \texttt{bad\_A} means violation of \texttt{A},
\item
  \texttt{bad\_G} means violation of \texttt{G}.
\end{itemize}

These are exactly the artifacts consumed by the logical-product phase in
Section 6.

\subsubsection{5.3 Example (toggle)}\label{example-toggle}

For \texttt{toggle}, an obvious choice is:

\begin{itemize}
\tightlist
\item
  assumption \(A=true\);
\item
  guarantee
  \(G=(y=0)\land \mathbf{G}(y=0\Rightarrow X(y=1))\land \mathbf{G}(y=1\Rightarrow X(y=0))\).
\end{itemize}

The resulting safety automata for these formulas are shown below. For
this \(A=true\) case, a complete view may include an explicit sink
\(bad_A\), but it is unreachable and omitted here for readability.

\begin{center}
\includegraphics[width=0.38\linewidth]{spec/figures/toggle_assume.pdf}
\end{center}

\begin{center}
\includegraphics[width=0.72\linewidth]{spec/figures/toggle_guarantee.pdf}
\end{center}

\subsection{6. Logical Product and
Reachability}\label{logical-product-and-reachability}

We work on product triples
\((s,q_A,q_G)\in\mathcal S\times Q_A\times Q_G\), where \(s\) is the
program control state, \(q_A\) is the residual assumption state, and
\(q_G\) is the residual guarantee state.

Initialization is \((s_0,q^A_0,q^G_0)\).

Logical successor relation: \[
(s,q_A,q_G)\leadsto (s',q'_A,q'_G)
\] iff there exist a program transition
\(s\xrightarrow{\gamma/\rho}s'\), an automaton edge
\(q_A\xrightarrow{\alpha_A}q'_A\), and an automaton edge
\(q_G\xrightarrow{\alpha_G}q'_G\), such that \[
\mathrm{SAT}\bigl(\gamma\land \llbracket\alpha_A\rrbracket\land \llbracket\alpha_G\rrbracket\bigr).
\]

The product-compatibility check uses guards only. The update relation
\(\rho\) is handled later when generating local obligations.

Reachability is the least fixpoint: \[
Reach_0=\{(s_0,q^A_0,q^G_0)\},
\] \[
Reach_{k+1}=Reach_k\cup
\left\{
x'\ \middle|\
x=(s,q_A,q_G)\in Reach_k,\ x\leadsto x'=(s',q'_A,q'_G),\ q_G\neq bad_G,\ q'_G\neq bad_G
\right\},
\] \[
Reach=\bigcup_{k\in\mathbb N} Reach_k.
\]

So triples in \(bad_G\) are not propagated in \(Reach\). However,
candidate moves from reached triples toward \(q'_G=bad_G\) are still
enumerated later in Section 7 to generate \texttt{GenObl} formulas.

Implementation correspondence: this phase is a BFS over product triples,
where each candidate successor is validated by one SAT check of guard
compatibility.

\subsubsection{6.1 Examples (toggle and delay
products)}\label{examples-toggle-and-delay-products}

Example (\texttt{delay\_int} product):

Using \(N_{delay}\) (Section 3.1.3), take the contract pair \[
A=true,\qquad G=\mathbf{G}(X(y)=x)
\] which captures one-tick delay from the second tick onward. A
corresponding safety-automaton shape for \(G\) uses a tracking state
after initialization and a bad sink.

At product level, one representative reachable prefix is: \[
(Init,ok_A,q^G_{init})\leadsto(Run,ok_A,q^G_{run}),
\] then in \texttt{Run}: \[
(Run,ok_A,q^G_{run})\leadsto(Run,ok_A,q^G_{run})\leadsto\cdots
\] as long as step guards remain compatible with the delay condition.

An abstract reachable-product sketch is:

\begin{center}
\includegraphics[width=0.70\linewidth]{spec/figures/delay_product.pdf}
\end{center}

Example (\texttt{toggle} product):

Using \(N_{toggle}\) (Section 3.1.3) and the automata of Section 5.3: \[
(Init,ok_A,q_0)\leadsto(Run,ok_A,q_1)
\] through \(Init\xrightarrow{true/y'=0}Run\),
\(ok_A\xrightarrow{true}ok_A\), and \(q_0\xrightarrow{(y=0)}q_1\).

Then the \texttt{Run\ -\textgreater{}\ Run} transition
\(Run\xrightarrow{true/y'=1-y}Run\) alternates the guarantee state: \[
(Run,ok_A,q_1)\leadsto(Run,ok_A,q_0)\leadsto(Run,ok_A,q_1)\leadsto\cdots
\]

Automaton-style rendering of the reachable product:

\begin{center}
\includegraphics[width=0.72\linewidth]{spec/figures/toggle_product.pdf}
\end{center}

No reachable triple contains \(bad_G\), which matches the expected
correctness of \texttt{toggle}.

Why no \texttt{bad} state appears in these product views: Section 6 only
shows the reachable safe product, i.e., triples kept in \(Reach\). By
construction, triples with \(q_G=bad_G\) (and successors with
\(q'_G=bad_G\)) are not propagated in this reachability set. Potential
moves toward \(bad_G\) are still considered later in Section 7 to
generate \texttt{GenObl} formulas.

\subsection{7. Generating Local
Obligations}\label{generating-local-obligations}

\subsubsection{\texorpdfstring{7.1 Formal Definition (\texttt{GenHyp}
then
\texttt{GenObl})}{7.1 Formal Definition (GenHyp then GenObl)}}\label{formal-definition-genhyp-then-genobl}

Let \(Reach \subseteq \mathcal S \times Q_A \times Q_G\) be the reached
set.

For a fixed program transition \(t=(s,\gamma,\rho,s')\in\mathcal T\), we
enumerate compatible combinations starting from reached source triples
\((s,q_A,q_G)\in Reach\), with outgoing automaton edges from \(q_A\) and
\(q_G\).

\texttt{Compat} (locally SAT-compatible combinations):

\[
\begin{aligned}
Compat_t=\{&(q_A,q_G,q'_A,q'_G,\alpha_A,\alpha_G)\mid\\
&(s,q_A,q_G)\in Reach,\\
&(q_A,\alpha_A,q'_A)\in\Delta_A,\ (q_G,\alpha_G,q'_G)\in\Delta_G,\\
&\mathrm{SAT}(\gamma\land\llbracket\alpha_A\rrbracket\land\llbracket\alpha_G\rrbracket)\}.
\end{aligned}
\]

The SAT filter is guard-only:
\(\gamma\land\llbracket\alpha_A\rrbracket\land\llbracket\alpha_G\rrbracket\),
interpreted on the current finite history. The update relation \(\rho\)
is used later in the deductive obligations attached to \(t\), not in
this compatibility test.

For readability, for any
\(c=(q_A,q_G,q'_A,q'_G,\alpha_A,\alpha_G)\in Compat_t\), define: \[
\psi_t(c)=\gamma\land\llbracket\alpha_A\rrbracket\land\llbracket\alpha_G\rrbracket,
\] \[
safe(c)\iff (q'_A\neq bad_A \land q'_G\neq bad_G),\qquad
badG(c)\iff (q'_A\neq bad_A \land q'_G=bad_G).
\]

\texttt{GenHyp} (disjunction of safe incoming contexts for \(t\)): keep
combinations in \(Compat_t\) whose successors stay outside \(bad_A\) and
\(bad_G\), then build:

\[
Hyp_t=
\bigvee
\left\{
\psi_t(c)
\mid
c\in Compat_t,\ safe(c)
\right\}.
\]

If the set is empty, no hypothesis is added. Otherwise \(Hyp_t\) is
added to \texttt{requires(t)}.

\texttt{GenObl} (forbid bad-\(G\) successors under active assumptions):
keep combinations in \(Compat_t\) with \(q'_G=bad_G\) and
\(q'_A\neq bad_A\), then generate:

\[
Obl_t=
\left\{
\neg\psi_t(c)
\mid
c\in Compat_t,\ badG(c)
\right\}.
\]

All formulas in \(Obl_t\) are added to \texttt{ensures(t)}.

The condition \(q'_A\neq bad_A\) enforces assume/guarantee discipline:
once assumptions are violated, no guarantee-preservation obligation is
generated.

Generated annotations for \(t\): \[
(requires(t),ensures(t))\;\supseteq\;(Hyp_t,Obl_t).
\]

Global phase:

\[
GenHypObl(N,\mathcal A,\mathcal G,Reach)=
\left(\{Hyp_t\}_t,\{Obl_t\}_t\right).
\]

Abstractly, Section 7 computes a transition annotation map: \[
Ann_7:\mathcal T\to \mathcal P(\mathrm{HFO}(X))\times \mathcal P(\mathrm{HFO}(X)),
\] with \(Ann_7(t)=(Req_t^{(7)},Ens_t^{(7)})\), where \(Req_t^{(7)}\) is
inserted in \texttt{requires(t)} and \(Ens_t^{(7)}\) in
\texttt{ensures(t)}.

\subsubsection{7.2 Examples (delay and
toggle)}\label{examples-delay-and-toggle}

Example (\texttt{delay\_int}):

Using the delay product from Section 6.1, consider the two reached safe
contexts: \[
(Init,ok_A,q^G_{init}),\qquad (Run,ok_A,q^G_{run}).
\]

For \texttt{GenHyp}, the \texttt{Init\ -\textgreater{}\ Run} transition
gets one compatible safe combination, yielding: \[
Hyp_{init}\coloneqq \gamma_{init}\land true\land\llbracket\alpha_{init}\rrbracket,
\] where \(\alpha_{init}\) is the guarantee-edge label from
\(q^G_{init}\) to its safe successor.

For the loop transition \texttt{Run\ -\textgreater{}\ Run}, the
compatible safe combination yields: \[
Hyp_{delay}\coloneqq \gamma_{run}\land true\land\llbracket\alpha_{run}\rrbracket.
\] Since \(\gamma_{run}=true\) in \(N_{delay}\), this simplifies to: \[
Hyp_{delay}=\llbracket\alpha_{run}\rrbracket.
\]

For \texttt{GenObl}, bad-target combinations are those with
\(q'_G=bad_G\) and \(q'_A\neq bad_A\); each one produces an
\texttt{ensures} formula of the form: \[
\neg\bigl(\gamma\land true\land\llbracket\alpha_{bad}\rrbracket\bigr).
\] These obligations are attached to the corresponding transition and
rule out locally compatible moves toward \(bad_G\).

Example (\texttt{toggle}):

For \texttt{toggle}, the reachable product in Section 6.1 has only safe
alternating triples and no compatible bad-target combination.

Therefore:

\[
Obl_t=\varnothing
\]

for all transitions \(t\).

For \texttt{GenHyp}, the result is non-empty on the
\texttt{Run\ -\textgreater{}\ Run} transition. Using the two reachable
product contexts \((Run,ok_A,q_0)\) and \((Run,ok_A,q_1)\), the
compatible safe combinations induce:

\[
HypLoop \coloneqq (\gamma\land true\land(y=0)) \lor (\gamma\land true\land(y=1)).
\]

Since \(\gamma=true\) on \texttt{Run\ -\textgreater{}\ Run} for
\texttt{toggle}, this simplifies to:

\[
HypLoop=(y=0)\lor(y=1),
\]

which is added to \texttt{requires(Run\ -\textgreater{}\ Run)}. This
context fact is simplifiable, but it is intentionally kept in generated
annotations for debugging and traceability.

\subsection{\texorpdfstring{8. User-Provided State Invariants
(\texttt{Inv})}{8. User-Provided State Invariants (Inv)}}\label{user-provided-state-invariants-inv}

User invariants are proof annotations attached to control states through
\[
Inv:\mathcal S\to \mathrm{HFO}(X).
\]

In the implementation, \texttt{Inv} is built and injected by a dedicated
coherency pass.

Construction of \(Inv(s)\):

\begin{enumerate}
\def\labelenumi{\arabic{enumi}.}
\tightlist
\item
  collect declared state invariants already attached to state \(s\),
\item
  collect user \texttt{ensures} formulas targeting state \(s\),
\item
  build \(Inv(s)\) as their conjunction (after deduplication).
\end{enumerate}

Insertion on transitions \(t:s\to s'\):

\[
requires(t)\;\leftarrow\;requires(t)\cup\{Inv(s)\},
\] \[
ensures(t)\;\leftarrow\;ensures(t)\cup\{Shift(Inv(s'))\}.
\]

An initialization coherency goal is also added at \texttt{init\_state}:
\[
Shift(Inv(s_0)).
\]

\texttt{Shift} is defined abstractly by semantic correctness: for any
\(\phi\in\mathrm{HFO}(X)\) and any finite history
\(h\in\mathrm{Hist}(X)\), \[
h \models Shift(\phi)\iff next(h)\models \phi,
\] where \(next(h)\) is the one-step extended history.

This insertion is purely proof-oriented: it does not modify the
operational step semantics. In pure-automate mode, these state-relation
invariants are then reused to strengthen local transition contexts
during obligation generation.

Combined with Section 7, this yields the final per-transition proof
annotation: \[
Ann:\mathcal T\to \mathcal P(\mathrm{HFO}(X))\times \mathcal P(\mathrm{HFO}(X)),
\] where for each \(t\), \(Ann(t)=(Req_t,Ens_t)\) is the pair of formula
sets ultimately attached to \texttt{requires(t)} and
\texttt{ensures(t)}.

\subsection{9. Transition-Local Deductive
Verification}\label{transition-local-deductive-verification}

At this stage, each transition \(t\in\mathcal T\) carries two sets of
formulas: \[
Req_t\subseteq \mathrm{HFO}(X),\qquad Ens_t\subseteq \mathrm{HFO}(X),
\] coming from Sections 7 and 8 (plus user contracts when present).

In the implementation, verification conditions are generated as guarded
local implications. For a transition \(t=(s,\gamma,\rho,s')\), define
its activation condition: \[
Active_t \;\coloneqq\; (state=s)\land \gamma.
\]

The generated obligations are abstractly: \[
\mathrm{Req\text{-}VC}(t):\quad Active_t \Rightarrow \bigwedge Req_t,
\] \[
\mathrm{Ens\text{-}VC}(t):\quad Active_t \land \rho \land \bigwedge Req_t \Rightarrow \bigwedge Ens_t.
\]

\texttt{Req-VC} comes from transition \texttt{requires}: if \(t\) is
active, its \texttt{requires} must hold.\\
\texttt{Ens-VC} comes from transition \texttt{ensures}: for an active
\(t\), under \(\rho\) and established \texttt{requires}, the
\texttt{ensures} must hold.

So guards are not a separate side condition: they are part of the
activation predicate that gates transition-local obligations.

Equivalent relational/inclusion view (Hoare-style with a relational
transition):

For \(t=(s,\gamma,\rho,s')\), define: \[
Pre_t(h)\coloneqq Active_t(h)\land \bigwedge Req_t(h),
\] \[
Post_t(h')\coloneqq \bigwedge Ens_t(h'),
\] \[
Rel_t(h,h')\coloneqq \rho(h,h').
\]

Then \texttt{Ens-VC} is equivalently: \[
\forall h,h'.\; Pre_t(h)\land Rel_t(h,h')\Rightarrow Post_t(h'),
\] or, as a set inclusion: \[
Rel_t[\{h\mid Pre_t(h)\}] \subseteq \{h'\mid Post_t(h')\}.
\]

This is the relational counterpart of weakest-precondition reasoning: no
imperative code block is needed, because the middle of the Hoare triple
is the transition relation itself.

This is Hoare-style in spirit, but operationally it is encoded as a
family of guarded implications that are composed in the global VC.

\subsection{10. Soundness Theorem (with Key
Lemmas)}\label{soundness-theorem-with-key-lemmas}

\subsubsection{10.1 Lemma 1 (Local generation
exactness)}\label{lemma-1-local-generation-exactness}

Every local combination \((t,q_A,q_G,q'_A,q'_G,\alpha_A,\alpha_G)\)
leading to \texttt{bad\_G} under \(q'_A \neq bad_A\) (and encountered
during reached-triple enumeration) produces an obligation that exactly
negates its guard conjunction.

\subsubsection{10.2 Lemma 2 (Local safety)}\label{lemma-2-local-safety}

If all local obligations of a transition \texttt{t} are valid, then no
concrete step executing \texttt{t} can realize a local combination
leading to \texttt{bad\_G} under \(q'_A \neq bad_A\).

Proof: immediate by contradiction, each forbidden combination is the
negation of one precise conjunction.

\subsubsection{\texorpdfstring{10.3 Lemma 3 (Preservation outside
\texttt{bad\_G})}{10.3 Lemma 3 (Preservation outside bad\_G)}}\label{lemma-3-preservation-outside-bad_g}

Assume all obligations are valid. For any concrete step in the trace and
any associated coherent logical triple, if we are outside
\texttt{bad\_G} before the step and outside \texttt{bad\_A} after the
step, then we remain outside \texttt{bad\_G} after the step.

Proof: apply Lemma 2 to the current step.

\subsubsection{10.4 Main Theorem}\label{main-theorem}

\paragraph{Statement}\label{statement}

Under the assumptions:

\begin{enumerate}
\def\labelenumi{\arabic{enumi}.}
\tightlist
\item
  correctness of LTL-to-automata compilation (\texttt{A}, \texttt{G}),
\item
  correctness of guard interpretation as first-order predicates over
  finite histories,
\item
  WP/backend correctness (all proved obligations are valid),
\item
  adequacy of local enumeration (all relevant combinations are covered),
\end{enumerate}

then:

\[
\forall \rho,\; \rho\!\upharpoonright_{I} \models A
\implies
\rho\!\upharpoonright_{I\cup O} \models G.
\]

\paragraph{Proof (outline)}\label{proof-outline}

\begin{itemize}
\tightlist
\item
  Let \(\rho\) be an execution such that
  \(\rho\!\upharpoonright_{I} \models A\).
\item
  By correctness of \(\mathcal A\), the \texttt{A} automaton never
  visits \texttt{bad\_A}.
\item
  Assume for contradiction that
  \(\rho\!\upharpoonright_{I\cup O} \not\models G\).
\item
  Then, by correctness of \(\mathcal G\), there exists a first index
  where the \texttt{G} state reaches \texttt{bad\_G}.
\item
  On the corresponding step, we have a local combination
  \((t,\alpha_A,\alpha_G)\) that leads to \texttt{bad\_G} while staying
  outside \texttt{bad\_A}.
\item
  This combination is covered by generation (Assumption 4), so its local
  obligation exists (Lemma 1), and is valid (Assumption 3), hence
  impossible (Lemma 2).
\item
  Contradiction.
\end{itemize}

Therefore \(\rho\!\upharpoonright_{I\cup O} \models G\).

\subsection{11. Implementation Mapping and Pass
Order}\label{implementation-mapping-and-pass-order}

From the implementation perspective, the pipeline is:

\begin{enumerate}
\def\labelenumi{\arabic{enumi}.}
\tightlist
\item
  parse/normalize source contracts and abstract node structure,
\item
  compute safety automata for \texttt{A} and \texttt{G},
\item
  build and explore the logical product over \((s,q_A,q_G)\),
\item
  run \texttt{GenHyp} to add incoming-context formulas to transition
  \texttt{requires},
\item
  run \texttt{GenObl} to add bad-target forbiddance formulas to
  transition \texttt{ensures},
\item
  discharge generated obligations with standard WP/SAT backends.
\end{enumerate}

Operational details:

\begin{itemize}
\tightlist
\item
  obligations are generated as local first-order formulas over finite
  histories,
\item
  compatibility checks use SAT on guard conjunctions,
\item
  the concrete implementation of \texttt{Shift} is realized with history
  operators (\texttt{now}/\texttt{pre\_k}) and checked against the
  abstract criterion given in Section 8,
\item
  generated formulas are simplified before insertion (including trivial
  ones, kept for diagnostics),
\item
  the concrete generated artifact is an annotated OBC+ node where each
  transition carries the computed \texttt{requires}/\texttt{ensures}
  pair,
\item
  verification is performed in pure-automate mode.
\end{itemize}

\subsection{12. Examples End-to-End}\label{examples-end-to-end}

The two running examples cover the full chain:

\begin{enumerate}
\def\labelenumi{\arabic{enumi}.}
\tightlist
\item
  \texttt{delay\_int}: abstract syntax (Section 3.1.3) and stream
  semantics (Section 4.2) show one-tick delay behavior.
\item
  \texttt{toggle}: contract automata (Section 5.3), reachable logical
  product (Section 6.1), and local obligation generation (Section 7.2)
  illustrate the end-to-end method.
\end{enumerate}

In both examples, the generated local formulas match the expected
temporal intent after simplification.

\subsection{13. Limits and Assumptions}\label{limits-and-assumptions}

This formalization relies on:

\begin{itemize}
\tightlist
\item
  deterministic program profile (\texttt{Det-Step},
  \texttt{Input-Frame}),
\item
  contracts in the safety-oriented fragment (positivity condition on
  \texttt{W}),
\item
  correctness of contract-to-automata construction,
\item
  adequacy of product reachability/enumeration for relevant
  combinations,
\item
  soundness of SAT/WP backends used to discharge obligations.
\end{itemize}

A practical caution is that systematically trivial generated formulas
(for example \(\neg(true\land true\land true)\)) often indicate a
construction issue, typically around bad-state propagation.

\end{document}
